\chapter{Hardware Heterogéneo: La Matriz de Despacho y el Contrato de Silicio}


El desarrollo asintótico de POLYDIM para cardinalidades \(D \ge 10^6\) impone un desafío monumental en términos de ejecución física. La presunción fundamental de la programación cognitiva moderna radica en que el agente de Inteligencia Artificial debe operar nativamente en espacios de alta dimensión \(S^{D-1}\), requiriendo un acoplamiento directo y sin fricciones con el hardware subyacente. Este capítulo formaliza el ``Contrato de Silicio'' (Silicon Contract), una política estricta de agnosticismio de hardware que prohíbe el uso de constantes mágicas y exige la interrogación dinámica de la topología del sistema. A través de la Matriz de Despacho, POLYDIM distribuye la carga computacional óptimamente entre procesadores x86-64, aceleradores GPU (NVIDIA CUDA, AMD ROCm), Google TPUs, arquitecturas de escala de oblea (Cerebras WSE-3) y computadoras cuánticas (QPUs).

\section{El Contrato de Silicio: Interrogar, No Asumir}


El Contrato de Silicio es un principio fundacional de la arquitectura POLYDIM que dictamina que el software jamás debe presuponer las especificaciones físicas de su entorno de ejecución. En sistemas de cómputo heterogéneo y topologías de enjambre de agentes, codificar rígidamente (hardcode) parámetros como el tamaño de la página de memoria, las líneas de caché, la latencia de interconexión o el ancho de los registros SIMD conduce indefectiblemente a fallos asintóticos catastróficos o cuellos de botella de rendimiento no portables.

\subsection{Diseño de la Clase HardwareProbe}
Para materializar este contrato, hemos diseñado la clase \texttt{HardwareProbe}. Esta abstracción actúa como el sismógrafo principal del Orquestador, encargada de perfilar la topología del CPU, la RAM disponible y la presencia de aceleradores dedicados (CUDA, ROCm, TPU) en tiempo de ejecución. 

\begin{verbatim}
class HardwareProbe:
    @staticmethod
    def detect_topology() -> Dict[str, Any]:
        topology = {
            "cpu_cores": os.cpu_count(),
            "ram_bytes": psutil.virtual_memory().total,
            "page_size": os.sysconf("SC_PAGE_SIZE") if hasattr(os, "sysconf") else 4096,
            "cuda_available": torch.cuda.is_available() if torch_installed else False,
            "rocm_available": torch.version.hip is not None if torch_installed else False,
            "tpu_available": check_tpu_presence(),
            "eps_float64": np.finfo(np.float64).eps,
            "arch_bits": 64 if sys.maxsize > 2**32 else 32
        }
        return topology
\end{verbatim}

\subsection{Anti-Hardcoding y Tolerancia de Máquina}
El uso de constantes mágicas como \texttt{1e-8} para la tolerancia al error (épsilon) está estrictamente prohibido. En su lugar, el algoritmo debe adaptarse a la precisión del tipo de dato, consultando \texttt{np.finfo(np.float64).eps}. Del mismo modo, el tamaño de la página (\texttt{os.sysconf('SC\_PAGE\_SIZE')}) es vital para alinear los tensores compartidos en memoria (PMTP Zero-Copy IPC) y evitar el fenómeno de \textit{false sharing} o la fragmentación en el Translation Lookaside Buffer (TLB) cuando \(D = 10^7\). El tamaño de la arquitectura se determina empíricamente consultando \texttt{sys.maxsize}.

\section{Intel/AMD x86-64: La Plataforma de Referencia}


La arquitectura x86-64 sigue siendo la plataforma de referencia ineludible y el ancla de estabilidad de POLYDIM. Garantiza el \textit{fallback} absoluto en cualquier nodo del enjambre.

\subsection{Paralelismo OpenMP y Afinidad de Hilos}
La implementación en C++ aprovecha la paralelización compartida mediante OpenMP (\texttt{\#pragma omp parallel for simd}). Se utiliza \texttt{omp\_get\_max\_threads()} para dimensionar estáticamente las regiones paralelas, asegurando que la afinidad de hilos (thread affinity) esté correctamente mapeada a la topología física, minimizando la migración de hilos entre núcleos, lo cual destruye la localidad espacial de los tensores.

\subsection{Vectorización AVX-512 y Jerarquía de Caché}
Los procesadores modernos soportan Advanced Vector Extensions (AVX-512), permitiendo el procesamiento simultáneo de 8 escalares de precisión doble (\texttt{double}) por registro SIMD. Mediante el uso juicioso del calificador \texttt{\_\_restrict\_\_}, indicamos al compilador (GCC 14) que los punteros no tienen solapamiento (aliasing), habilitando la autovectorización agresiva.

La jerarquía de caché (L1 64KB, L2 512KB, L3 32MB) dicta la estrategia de reducción. El algoritmo de suma compensada de Neumaier se divide en bloques (tiles) que caben perfectamente en L2, mitigando la latencia de DRAM. 
Bajo estas optimizaciones, el tiempo de ejecución (\textit{Benchmark}) para \(D=10^6\) alcanza los \textbf{35.06 ms}, exhibiendo una escalabilidad lineal estricta \(O(D)\). Para entornos de Alta Computación (HPC), se habilita el flag \texttt{POLYDIM\_ENABLE\_FTZ} (Flush-To-Zero) para subnormales, evitando penalizaciones de microcódigo en underflows.

\section{NVIDIA CUDA vía Triton: Aceleración Tensorial}


Para trascender la limitación del ancho de banda de memoria principal del CPU, POLYDIM implementa núcleos de GPU nativos utilizando Triton, un lenguaje incrustado en Python diseñado para escribir núcleos (kernels) de GPU de alta eficiencia sin descender al nivel de PTX o CUDA C.

\subsection{Arquitectura polydim\_triton\_kernel\_v764.py}
El núcleo se diseña con un \texttt{BLOCK\_SIZE} estático de 1024 elementos por bloque de hilos. Dado que las operaciones en la esfera unitaria \(S^{D-1}\) requieren normativas hiperdimensionales, se implementa una reducción en bloques para Neumaier en precisión FP64. 

\begin{itemize}
    \item \textbf{Carga con Máscaras:} Se utiliza \texttt{tl.load} con una máscara condicional para manejar con seguridad los bordes de tensores cuyas dimensiones \(D\) no son múltiplos exactos del \texttt{BLOCK\_SIZE}.
    \item \textbf{Reducción Atómica:} \texttt{tl.atomic\_add} permite la reducción cruzada de bloques de forma segura en memoria global (VRAM).
\end{itemize}

El tensor reside perpetuamente en la memoria de video (VRAM) mediante Direct DMA, eliminando el latido y el roundtrip al CPU. El perfil de latencia exhibe un costo fijo de lanzamiento del núcleo (kernel launch) de \(\approx 50 \mu s\). Para \(D=10^6\), el benchmark alcanza \textbf{4.12 ms}, representando una aceleración de 8.5x frente al CPU x86-64.

\section{AMD ROCm / HIP: Computación Heterogénea}


La interfaz HIP (Heterogeneous-compute Interface for Portability) permite que el núcleo matemático abstracto de POLYDIM sea compilado en el formato binario HSACO (Hardware Abstraction Shader Compiler Object) nativo de AMD. 

La principal divergencia arquitectónica respecto a NVIDIA es la unidad mínima de ejecución: AMD emplea un Wave64 de 64 hilos, en contraste con el Warp de 32 hilos de NVIDIA. Esto requiere un ajuste dinámico en el \textit{loop unrolling} y en el despachador de bloques de memoria compartida dentro del kernel.
A fecha de 2026, la política oficial de POLYDIM restringe el uso de HIP a entornos Linux/Cloud, dada la histórica inestabilidad del soporte ROCm en sistemas Windows nativos.
Benchmark a \(D=10^6\): \textbf{5.80 ms}.

\section{Google Cloud TPU: Álgebra Lineal Acelerada}


Las Unidades de Procesamiento Tensorial (TPU) de Google, específicamente la generación v3-8, ofrecen una ventaja algorítmica fundamental mediante su arquitectura de arreglo sistólico, que ejecuta multiplicaciones matriciales densas en tiempo \(O(n)\) aprovechando un hardware espacial de tamaño \(O(n^2)\).

La integración se realiza a través de XLA (Accelerated Linear Algebra) mediante el pipeline de compilación \textit{Just-In-Time} (JIT) de JAX. Dado que la TPU opera nativamente en formato Bfloat16 (BF16), la precisión FP64 de POLYDIM se emula por hardware/software, introduciendo un ligero overhead. Sin embargo, para proyecciones de Johnson-Lindenstrauss utilizando la Transformada Rápida de Walsh-Hadamard (FWHT), el procesamiento paralelo masivo es insuperable.
Benchmark a \(D=10^6\) en TPU v3-8: \textbf{2.30 ms}.

\section{Cerebras WSE-3 (CS-3): El Sueño POLYDIM}


El Wafer-Scale Engine 3 (WSE-3) de Cerebras System redefine los límites físicos del hardware de IA. Con 4 billones de transistores, 900,000 núcleos de IA y 44 GB de memoria SRAM directamente en el chip (on-chip), representa el Santo Grial para POLYDIM: \textbf{Cero Cuello de Botella de Memoria DRAM}.

En las arquitecturas de von Neumann tradicionales (CPU/GPU), mover los datos a la memoria externa consume órdenes de magnitud más energía y tiempo que la operación aritmética real. En WSE-3, toda la memoria es local, y la interconexión de malla 2D permite una comunicación directa núcleo-a-núcleo sin pasar por buses periféricos. 
Implementado mediante el Cerebras Software Language (CSL), el benchmark pulveriza los registros con \textbf{0.85 ms} a \(D=10^6\), convirtiéndolo en el nodo más rápido de la Matriz de Despacho.

\section{QPU (Unidades de Procesamiento Cuántico)}


El horizonte terminal del despliegue geométrico de POLYDIM son las computadoras cuánticas basadas en compuertas universales (IBM Quantum, Rigetti, IonQ, Google Sycamore). El conjunto universal de compuertas Clifford+T permite compilar rotaciones en \(\text{SU}_q(2)\) a cualquier hardware QPU compatible.

La principal restricción actual es la coherencia cuántica, dictada por los tiempos límite \(T_1\) (relajación de amplitud) y \(T_2\) (desfase). Además, modelar \(D=10^6\) estados requiere, como mínimo teórico estricto, 20 qubits completamente lógicos y entrelazados (\(2^{20} > 10^6\)). Para mitigar el ruido durante la deformación tensorial a través del enjambre, se aplica el protocolo de \textit{Clifford Twirling}.

\section{La Matriz de Despacho Asintótica}


A continuación, la tabla maestra de enrutamiento asintótico consolidada:


\centering
\begin{tabular}{|l|l|l|l|l|}
\hline
\textbf{Plataforma} & \textbf{Backend} & \textbf{Precisión Nativa} & \textbf{Ventaja Arquitectónica} & \textbf{Latencia (\(D=10^6\))} \\ \hline
Intel/AMD x86-64    & OpenMP + AVX-512 & FP64                      & Compatibilidad universal        & 35.06 ms                      \\ \hline
NVIDIA CUDA         & Triton GPU JIT   & FP64/FP32                 & Ancho de banda HBM2e            & 4.12 ms                       \\ \hline
AMD ROCm / HIP      & HSACO (Linux)    & FP64/FP32                 & Wave64 Compute                  & 5.80 ms                       \\ \hline
Google Cloud TPU    & XLA (JAX)        & BF16 (FP64 emul.)         & Arreglo Sistólico O(n)          & 2.30 ms                       \\ \hline
Cerebras WSE-3      & CSL              & FP16/FP32                 & Zero DRAM, 44GB SRAM en chip    & 0.85 ms                       \\ \hline
QPU (Gate-based)    & Clifford+T       & Amplitud Cuántica         & Superposición de estados        & Lim. Coherencia (T1/T2)       \\ \hline
\end{tabular}


